% !TeX root = ../../Book3.tex
% 纲要标题：## 第6章　Hilbert 零点定理、正规化与导子理想
% 纲要位置：工作记录/计划纲要.md，第 2062 行。

\chapter{Hilbert 零点定理、正规化与导子理想}\label{b3:ch:06}
\BookChapterNumber{b3:ch:06}

\begin{chapterabstract}
零点定理把代数闭域上的多项式零点集与根理想联系起来，正规化则补入分式域或全分式环中的整元素。本章先由 Zariski 引理证明弱零点定理，再用 Rabinowitsch 技巧得到强形式，随后建立正规化的对象，完整证明任意域上有限型代数的正规化有限性。尖点、节点及导子理想的计算，用来区分有限性、双有理性、点集纤维与局部环的变化。
\end{chapterabstract}

\begin{learninggoals}
\begin{itemize}
\item 区分 \(L=k(\alpha_1,\ldots,\alpha_r)\) 的域生成与 \(L=k[\alpha_1,\ldots,\alpha_r]\) 的代数生成，判断极大理想的剩余域。
\item 从零点条件得到根理想成员关系，并写清 Rabinowitsch 技巧中的分母消去。
\item 对整环、约化环和非约化环采用明确的正规化约定。
\item 分开处理可分迹论证与正特征不可分扩张，指出正规化有限性证明中需要添加的有限组系数。
\item 计算曲线正规化、导子及商模，避免把点数、环同构和平坦性混为一谈。
\end{itemize}
\end{learninggoals}

在代数闭域 \(k\) 上，方程 \(x^2=0\) 与 \(x=0\) 有完全相同的零点；然而 \(k[x]/(x^2)\) 中的 \(\bar x\) 仍非零，且平方为零。只看取值点会忘记这个幂零元素。零点定理所对应的不是原理想本身，而是它的根；要把这一点推广到多个变量，先须知道极大理想是否都来自实际的取值点。

\ProseParagraphBreak

另一种失真发生在一个环已没有幂零元时：尖点的坐标环 \(k[t^2,t^3]\) 仍缺少分式域中的元素 \(t\)，虽然它满足首一方程。把所有这样的整元素补入，会得到正规化。它既保持分式域，又可能改变局部环和一点之上的纤维。因此零点集、坐标环与正规环是三个需要分别比较的对象。

\ProseParagraphBreak

本章前半使用 Noether 正规化、极大理想与根理想；后半还用到有限生成模的 Noether 性，以及第一卷的可分扩张、迹配对和纯不可分扩张。零点定理已足以支持\secref{b3:sec:0701}与\secref{b3:sec:0702}的坐标环、正则函数和态射；\secref{b3:sec:0604}的两条曲线与\secref{b3:sec:0606}的导子计算还使用具体的正规化。两条曲线的有限性可由显式模生成式直接核验。\cref{b3:prop:reduced-total-quotient-product}的分支分解则用于\secref{b3:sec:0704}的交叉直线与\chapref{b3:ch:20}的正规性判据。

\ProseParagraphBreak
Matsumura 的《Commutative Ring Theory》\cite[§5, Theorems 5.3--5.4, pp. 33--34]{Matsumura1989CommutativeRingTheory}可用于对照零点定理和 Rabinowitsch 技巧；正规化有限性的论证见该书\cite[§33, Integral closure of a Noetherian integral domain, pp. 261--263]{Matsumura1989CommutativeRingTheory}。本章把其中的不可分扩张处理和有限分支推广单独展开，逐项说明所需条件。

% !TeX root = ../../Book3.tex
% 纲要标题：### 6.1 Zariski 引理
% 纲要位置：工作记录/计划纲要.md，第 2064 行。
% 纲要细目（写作范围）：
% * 作为域的有限生成代数；域生成与代数生成、有限分母的限制
% * 极大理想的剩余域
% * 与代数闭域假设的联系；闭点扩域后的分裂

\section{Zariski 引理}\label{b3:sec:0601}

把一个有限型代数模掉极大理想，得到的剩余域依然由有限个元素作为代数生成。域的倒数运算与代数生成元的有限性相互约束，结果是这种剩余域只能是基域的有限代数扩张。Noether 正规化使这一事实有一个简洁证明。

\subsection{作为域的有限生成代数}\label{b3:subsec:0601:01}

设 \(k\) 为任意域，\(t\) 在 \(k\) 上超越。\(k(t)\) 虽由一个元素 \(t\) 作为域生成，却必须包含 \(k[t]\) 中每个非零多项式的倒数。有限个有理函数的分母乘积若为 \(D\ne0\)，则它们生成的 \(k\) 代数包含于 \(k[t,1/D]\)，其元素的分母都可写成 \(D\) 的幂。取非恒定多项式 \(1+tD\) 的一个不可约因子 \(h\)，则 \(h\nmid D\)，所以 \(h^{-1}\) 不在这个子环中：否则存在 \(P\in k[t]\) 和整数 \(m\ge0\)，使 \(D^m=hP\)，与 \(h\nmid D\) 矛盾。这说明有限个代数生成元不能提供这个超越方向上所需的全部倒数；下面的引理把这种限制推广到任意有限型域代数。

\begin{lemma}[Zariski]\label{b3:lem:zariski}
设 \(k\) 为域，\(L/k\) 为域扩张，且 \(L\) 作为 \(k\) 代数有限生成，则 \([L:k]<\infty\)。这一结论称为\term[zariski-yinli]{Zariski 引理}[Zariski's lemma]。
\end{lemma}
\begin{proof}
由\cref{b3:thm:noether-normalization}，存在多项式子环 \(R=k[y_1,\ldots,y_d]\subseteq L\)，使 \(L\) 为有限 \(R\) 代数，因而在 \(R\) 上整。因为 \(L\) 是域，\cref{b3:lem:integral-field-criterion}迫使 \(R\) 也是域。若 \(d>0\)，变量 \(y_1\) 在多项式环中非零且不可逆，矛盾。因此 \(d=0\)，\(R=k\)，而正规化结论正说明 \(L\) 是有限生成的 \(k\) 模，即有限维 \(k\) 向量空间。
\end{proof}

反过来，有限代数扩张选一个有限的向量空间基即可作为代数生成集，所以引理刻画了哪些域是有限型 \(k\) 代数。

\subsection{极大理想的剩余域}\label{b3:subsec:0601:02}

\begin{corollary}\label{b3:cor:finite-type-residue-fields}
设 \(k\) 为域，\(A\) 为交换含幺有限型 \(k\) 代数。对 \(A\) 的任一极大理想 \(\mathfrak m\)，剩余域 \(A/\mathfrak m\) 是 \(k\) 的有限扩张。对 \(A\) 的一般素理想 \(\mathfrak p\)，剩余域 \(\kappa(\mathfrak p)=\operatorname{Frac}(A/\mathfrak p)\) 则只必然是有限生成的域扩张。
\end{corollary}
\begin{proof}
\(A/\mathfrak m\) 是域，同时是 \(A\) 的有限型代数商，应用\cref{b3:lem:zariski}即可。对一般素理想，\(A/\mathfrak p\) 的有限组代数生成元也作为域生成元生成其分式域，但此时已经允许把任意非零多项式倒置，不能继续应用 Zariski 引理。
\end{proof}

例如 \(k[x,y]\) 的极大理想若取 \((x,y)\)，剩余域为 \(k\)；素理想 \((x)\) 的剩余域却是 \(k(y)\)。闭点和一般点的差别，在这里反映为剩余域的代数性与超越性。

\subsection{代数闭域假设的作用}\label{b3:subsec:0601:03}

若 \(k\) 代数闭，它没有非平凡有限代数扩张：有限扩张中的每个元素有 \(k\) 上的不可约最小多项式，而代数闭性迫使该多项式为一次。因此\cref{b3:cor:finite-type-residue-fields}给出
\[
A/\mathfrak m=k
\]
的典范 \(k\) 代数同构。这里的等号意为结构映射 \(k\to A/\mathfrak m\) 为同构，而不是选定某个任意的抽象域同构。

\ProseParagraphBreak
基域不代数闭时，有限剩余域扩张必须保留。例如 \(\mathbb R[x]/(x^2+1)\cong\mathbb C\)，对应一个实多项式环的极大理想，却不对应实数轴上的点。把系数域扩大到 \(\mathbb C\) 后，\(x^2+1=(x-i)(x+i)\)，两因子互素，中国剩余定理给出
\[
\mathbb C[x]/(x^2+1)\cong\mathbb C\times\mathbb C.
\]
原来的一个闭点由此分裂成 \(i,-i\) 两个复坐标点。另一方面，\(\mathbb F_q[x]\) 的次数 \(r\) 的首一不可约多项式给出剩余域 \(\mathbb F_{q^r}\)。下一节正是把代数闭性用在排除这些额外的剩余域上。

% !TeX root = ../../Book3.tex
% 纲要标题：### 6.2 弱零点定理
% 纲要位置：工作记录/计划纲要.md，第 2070 行。
% 纲要细目（写作范围）：
% * 代数闭域上的极大理想；多变量求值映射的核
% * 公共零点的存在判据；坐标点与素谱点的区别
% * 非代数闭域的对照

\section{弱零点定理}\label{b3:sec:0602}

一个坐标点可以由所有在该点取零的多项式描述。设 \(k\) 为任意域，\(n\ge0\)，\(a=(a_1,\ldots,a_n)\in k^n\)，考虑取值映射 \(\operatorname{ev}_a:k[x_1,\ldots,x_n]\to k\)。逐次将 \(x_1,\ldots,x_n\) 替换成 \(a_1,\ldots,a_n\)，每次作差都是关于 \(x_i\) 的多项式在 \(x_i\) 与 \(a_i\) 处的取值之差，因而被 \(x_i-a_i\) 整除。把这些差相加，对任意 \(f\) 得到
\[
f(x_1,\ldots,x_n)-f(a_1,\ldots,a_n)
=\sum_{i=1}^n(x_i-a_i)g_i(x_1,\ldots,x_n),
\qquad g_i\in k[x_1,\ldots,x_n].
\]
因此 \(\ker\operatorname{ev}_a=(x_1-a_1,\ldots,x_n-a_n)\)。取值映射在常数上是恒等映射，因而满射到域 \(k\)，所以这个核是极大理想。代数闭域上的弱零点定理将说明，每个极大理想都以这种方式产生。

\subsection{代数闭域上的极大理想}\label{b3:subsec:0602:01}

\begin{theorem}[Hilbert 弱零点定理]\label{b3:thm:weak-nullstellensatz}
设 \(k\) 是代数闭域，\(n\ge0\)。则 \(k[x_1,\ldots,x_n]\) 的每个极大理想唯一写成
\[
\mathfrak m_a=(x_1-a_1,\ldots,x_n-a_n),\qquad a\in k^n.
\]
这称为\term[ruo-lingdian-dingli]{弱零点定理}[weak Nullstellensatz]。
\end{theorem}
\begin{proof}
对极大理想 \(\mathfrak m\)，由\cref{b3:cor:finite-type-residue-fields}及代数闭性，结构映射 \(k\to k[x_1,\ldots,x_n]/\mathfrak m\) 是同构。令 \(a_i\in k\) 为 \(x_i\) 在商域中的像，则 \(x_i-a_i\in\mathfrak m\)，故 \(\mathfrak m_a\subseteq\mathfrak m\)。二者都极大，只能相等。若 \(\mathfrak m_a=\mathfrak m_b\)，将 \(x_i-a_i\) 在 \(b\) 处取值，得到 \(b_i=a_i\)，所以点唯一。
\end{proof}

\subsection{公共零点的存在判据}\label{b3:subsec:0602:02}

对理想 \(I\subseteq k[x_1,\ldots,x_n]\)，记
\[
Z_k(I)=\bigl\{\,a\in k^n\mid f(a)=0\text{ 对每个 }f\in I\,\bigr\}.
\]
\symidx{\(Z_k(I)\)：理想 \(I\) 在 \(k^n\) 中的公共零点集}
这与素谱中的 \(V(I)\) 需要区别：\(Z_k(I)\) 只记录 \(k\) 坐标点，而 \(V(I)\) 包含所有含 \(I\) 的素理想。若 \(k\) 代数闭，弱零点定理把 \(Z_k(I)\) 中的坐标点与 \(V(I)\) 中的闭点对应起来。例如取 \(I=(0)\subset k[x,y]\)，则 \(Z_k(0)=k^2\)，而 \(V(0)=\operatorname{Spec}k[x,y]\) 还包含 \((0),(x),(y-x^2)\) 等非闭点；它们的商环分别为 \(k[x,y],k[y],k[x]\)，都是整环却不是域，所以这些素理想都不极大。这些非闭点即使在 \(k\) 代数闭时也仍然存在。

\begin{corollary}\label{b3:cor:proper-ideal-common-zero}
设 \(k\) 为代数闭域，\(n\ge0\)，\(I\) 为 \(k[x_1,\ldots,x_n]\) 的理想。记 \(Z_k(I)=\{a\in k^n:f(a)=0\text{ 对所有 }f\in I\}\)。则 \(I\) 是真理想，当且仅当 \(Z_k(I)\ne\varnothing\)。特别地，给定整数 \(r\ge1\) 及 \(k[x_1,\ldots,x_n]\) 中的多项式 \(f_1,\ldots,f_r\)，它们没有 \(k^n\) 中的公共零点，当且仅当存在 \(g_1,\ldots,g_r\in k[x_1,\ldots,x_n]\) 满足
\[
g_1f_1+\cdots+g_rf_r=1.
\]
\end{corollary}
\begin{proof}
若 \(I\) 是真理想，可先取极大理想 \(\mathfrak m\supseteq I\)。\cref{b3:thm:weak-nullstellensatz}把这个纯代数的选择转为一个坐标点：\(\mathfrak m=\mathfrak m_a\)，其中 \(a\in k^n\)。于是每个 \(f\in I\) 都属于 \(\ker\operatorname{ev}_a\)，即 \(f(a)=0\)，所以 \(a\) 是公共零点。反过来，若 \(a\) 为公共零点，则 \(I\subseteq\ker\operatorname{ev}_a\)，所以 \(1\notin I\)。最后对 \(I=(f_1,\ldots,f_r)\) 使用理想生成的定义即可。
\end{proof}

例如 \(xy-1\) 和 \(x\) 没有公共零点，等式
\(yx-(xy-1)=1\) 就是可以直接核验的证据。零点定理保证任何无公共零点的有限系统都有这样的代数组合；它本身没有给出这些系数的次数上界。第一卷附录F.3的 Buchberger 算法及定理F.3.3给出一种计算办法：系数域支持精确四则运算和零判定，且选定的全局单项式序能够有效比较时，可从有限生成元计算 Gröbner 基；保存各步的生成元表示，再把 \(1\) 的约化过程回代，就能得到上述理想成员证据。

\subsection{非代数闭域上的对照}\label{b3:subsec:0602:03}

\((x^2+1)\subseteq\mathbb R[x]\) 是真理想，却没有实零点。有限域上的差别甚至会出现在零理想的零点理想中：多项式 \(x^q-x\) 在全部 \(\mathbb F_q\) 元素上取零，本身却不是零多项式。这些例子分别说明公共零点存在性及从点集恢复多项式关系，都依赖所使用的基域。

\ProseParagraphBreak
不过任何域上的真理想，在代数闭包中都有公共零点。对真理想 \(I\subseteq k[x_1,\ldots,x_n]\)，选择极大理想 \(\mathfrak m\supseteq I\)，其剩余域 \(L\) 由\cref{b3:cor:finite-type-residue-fields}是有限代数扩张。把 \(L\) 嵌入一个固定代数闭包 \(\overline k\)，变量的像给出所求点。这里有限扩域嵌入代数闭包，可按有限组代数生成元逐次选择最小多项式的根完成；不要求 \(L/k\) 可分。

\ProseParagraphBreak
若只研究实点或有限域上的有理点，就应当保留相应基域的额外关系，而不能把代数闭域版本原样照用。本卷后续的实代数补充也将据此区分不同的零点定理。

% !TeX root = ../../Book3.tex
% 纲要标题：### 6.3 强零点定理
% 纲要位置：工作记录/计划纲要.md，第 2076 行。
% 纲要细目（写作范围）：
% * 根理想与零点集的双向对应
% * Rabinowitsch 技巧；任意交换含幺环上的增广理想判据
% * 几何对象与代数对象反变关系的起点；普通坐标环的约化性

\section{强零点定理}\label{b3:sec:0603}

弱形式只判断方程组是否有解。强形式进一步回答：已知全部解之后，究竟能恢复哪些多项式关系？一个多项式与它的正整数次幂有相同零点，因此期望恢复原理想通常过强；正确对象是根理想。本节除特别说明外，\(k\) 始终为代数闭域。

\subsection{根理想与零点集的对应}\label{b3:subsec:0603:01}

设 \(k\) 为域，\(n\ge0\)，\(X\subseteq k^n\)。定义 \(X\) 的\term[lingdian-lixiang]{零点理想}[vanishing ideal]
\[
I(X)=\bigl\{\,f\in k[x_1,\ldots,x_n]\mid f(a)=0\text{ 对每个 }a\in X\,\bigr\}.
\]
这是根理想，因为 \(f(a)^r=0\) 在域中推出 \(f(a)=0\)。空集的零点理想按全称命题约定为整个多项式环。
\symidx{\(I(X)\)：点集 \(X\) 上恒为零的多项式组成的理想}

\begin{theorem}[Hilbert 强零点定理]\label{b3:thm:strong-nullstellensatz}
设 \(k\) 为代数闭域，\(n\ge0\)，\(J\) 为 \(k[x_1,\ldots,x_n]\) 的理想。记 \(Z_k(J)=\{a\in k^n:f(a)=0\text{ 对所有 }f\in J\}\)；对 \(X\subseteq k^n\)，记 \(I(X)=\{f\in k[x_1,\ldots,x_n]:f(a)=0\text{ 对所有 }a\in X\}\)。则
\[
I\bigl(Z_k(J)\bigr)=\sqrt J.
\]
这称为\term[qiang-lingdian-dingli]{强零点定理}[strong Nullstellensatz]。因此根理想与形如 \(Z_k(J)\) 的代数集，分别经 \(Z_k\) 与 \(I\) 构成反向包含的一一对应。
\end{theorem}

例如 \((x^2,y)\) 和 \((x,y)\) 在 \(k^2\) 中都只给出原点；从这个点集恢复的是根理想 \((x,y)\)。商环 \(k[x,y]/(x^2,y)\) 中保留的非零幂零元，无法由普通点集上的取值看到。下一小节证明等式，再核对对应关系。

\subsection{Rabinowitsch 技巧}\label{b3:subsec:0603:02}

把“不等于零”的条件写成“存在倒数”，是从弱形式过渡到强形式的关键。对 \(f\in k[x_1,\ldots,x_n]\)，再加变量 \(T\)，方程 \(Tf-1=0\) 正是要求 \(f\) 的值可逆。这个方法称为\term[rabinowitsch-jiqiao]{Rabinowitsch 技巧}[Rabinowitsch trick]。\footnote{原论文为《Zum Hilbertschen Nullstellensatz》，发表于 1930 年的《Mathematische Annalen》，见\cite{Rabinowitsch1930Nullstellensatz}。}

其中从增广理想回到根理想的推理，不依赖零点或代数闭性，而是下面的纯代数事实。

\begin{lemma}\label{b3:lem:rabinowitsch-equivalence}
设 \(R\) 为交换含幺环，\(I\subseteq R\) 为理想，\(f\in R\)，\(T\) 为不定元。则
\[
f\in\sqrt I
\quad\Longleftrightarrow\quad
IR[T]+(1-Tf)=R[T].
\]
\end{lemma}
\begin{proof}
若 \(f^N\in I\)、\(N\ge1\)，有限几何级数给出
\[
1=(1-Tf)\sum_{j=0}^{N-1}(Tf)^j+T^Nf^N,
\]
所以右侧理想含一。

反向，设右侧理想含一。由多项式环的泛性质，局部化映射 \(R\to R_f\) 延拓为环同态
\[
R[T]\longrightarrow R_f,\qquad T\longmapsto f^{-1}.
\]
在这个同态下，\(1-Tf\) 变成零，故单位理想的等式给出 \(1\in IR_f\)。将\cref{b3:prop:ideal-localization-saturation}用于乘法集 \(\{1,f,f^2,\ldots\}\)，得到 \(I\) 中含有某个 \(f^N\)。若 \(N=0\)，则 \(1\in I\)，可改取 \(N=1\)。因此 \(f\in\sqrt I\)。这里通过局部化的成员判据清分母，所以即使 \(R\) 有零因子，推理仍成立。
\end{proof}

\begin{proof}[\cref{b3:thm:strong-nullstellensatz}的证明]
若 \(f^r\in J\)，则在 \(Z_k(J)\) 的每个点都有 \(f(a)^r=0\)，从而 \(f(a)=0\)。因此 \(\sqrt J\subseteq I\bigl(Z_k(J)\bigr)\)。

反向，取在 \(Z_k(J)\) 上恒为零的多项式 \(f\)，在 \(k[x_1,\ldots,x_n,T]\) 中考虑
\[
J' = J\,k[x_1,\ldots,x_n,T]+(1-Tf).
\]
若 \((a,t)\) 为其零点，则 \(a\in Z_k(J)\)，所以 \(f(a)=0\)，但又要求 \(1-t\cdot f(a)=0\)，矛盾。由\cref{b3:cor:proper-ideal-common-zero}，无公共零点使 \(J'\) 成为单位理想；再由\cref{b3:lem:rabinowitsch-equivalence}，得到 \(f\in\sqrt J\)。弱零点定理在这里把几何上的无解转成单位理想，随后根理想的结论完全由纯代数引理给出。

两种运算都反向包含。刚证明对根理想 \(J\) 有 \(I\bigl(Z_k(J)\bigr)=J\)。对代数集 \(X=Z_k(J)\)，一方面 \(J\subseteq I(X)\) 给出 \(Z_k\bigl(I(X)\bigr)\subseteq X\)；另一方面每个 \(a\in X\) 按定义使 \(I(X)\) 中全部多项式为零，得到反向包含。于是两种运算在所述对象上互逆。
\end{proof}

\subsection{几何与代数的反变关系}\label{b3:subsec:0603:03}

\begin{definition}\label{b3:def:affine-coordinate-ring}
设 \(k\) 为代数闭域，\(n\ge0\)，\(X\subseteq k^n\) 为多项式方程组的公共零点集。定义商环
\[
k[X]=k[x_1,\ldots,x_n]/I(X)
\]
并称它为 \(X\) 的\term[zuobiao-huan]{坐标环}[coordinate ring]。其中一个元素就是多项式在 \(X\) 上给出的函数；两个多项式给出同一函数，当且仅当它们之差属于 \(I(X)\)。
\end{definition}
\symidx{\(k[X]\)：仿射代数集 \(X\) 的坐标环}

因为 \(I(X)\) 是根理想，普通代数集的坐标环总是约化环。若从理想 \(J\subseteq R=k[x_1,\ldots,x_n]\) 出发，先取零点集再取坐标环，强零点定理给出
\[
k[Z_k(J)]=R/\sqrt J.
\]
这未必是 \(R/J\)：例如前面的 \(J=(x^2,y)\) 得到原点的坐标环 \(k\)，而 \(R/J\) 中的非零幂零元 \(\bar x\) 已经丢失。

设 \(X,Y\subseteq k^n\) 为代数集。强零点定理说明，代数集的包含与零点理想的包含反向对应：
\[
X\subseteq Y\ \Longleftrightarrow\ I(Y)\subseteq I(X).
\]
当 \(X\subseteq Y\) 时，把 \(Y\) 上的多项式函数限制到 \(X\)，就得到满同态 \(k[Y]\rightarrow k[X]\)，其核是 \(I(X)/I(Y)\)。因此
\[
k[X]\cong k[Y]\big/\bigl(I(X)/I(Y)\bigr).
\]
例如 \(Y=k^2\)、\(X=Z_k(y)\) 时，限制函数把 \(k[x,y]\) 送到 \(k[x]\)，正是令 \(y=0\) 的商同态。代数集的包含由此表现为函数环的商，而不是子环。

\ProseParagraphBreak
推广到一般映射时，需要判断函数与映射复合后是否仍属于相应的函数环；在局部允许分母的情形，还要考察这些表达式在各点是否有定义。\secref{b3:sec:0701}与\secref{b3:sec:0702}以此研究代数集的拓扑、正则函数与态射。\secref{b3:sec:0604}的尖点、节点和\secref{b3:sec:0606}的导子计算，则进一步比较零点集与环本身保留的信息。

% !TeX root = ../../Book3.tex
% 纲要标题：### 6.4 正规化
% 纲要位置：工作记录/计划纲要.md，第 2082 行。
% 纲要细目（写作范围）：
% * 域上有限生成代数的正规化及其有限性问题：先处理整环，再处理约化环，非约化情形约定先取约化；定义与模型在本节建立，一般有限性定理在紧接的第6.5节陈述并证明
% * 正规化不有限的情形；缺失常数所造成的无限剩余域维数
% * 尖点和节点曲线的代数例子
% * 整闭包、有限性和双有理性的不同角色；本节的尖点、节点直接通过显式环包含核验，不依赖尚未证明的一般有限性定理

\section{正规化}\label{b3:sec:0604}

同一个有理函数域可能包含不同的坐标环。尖点的 \(k[t^2,t^3]\) 与直线的 \(k[t]\) 就有相同分式域，但前者漏掉了一个已经满足首一方程的函数 \(t\)。正规化把这样的整元素一并加入。它改变的是环；点集是否改变、扩张是否有限，需要分别核验。

\subsection{整环、约化环与正规化的定义}\label{b3:subsec:0604:01}

\begin{definition}\label{b3:def:normalization}
设 \(A\) 为交换含幺环。若 \(A\) 是整环，则称它在 \(\operatorname{Frac}(A)\) 中的整闭包 \(\widetilde A\) 为 \(A\) 的\term[zhengguihua]{正规化}[normalization]。若 \(A\) 是约化环，令 \(S\subseteq A\) 为其全部非零因子的集合，定义\term[quan-fenshi-huan]{全分式环}[total ring of fractions]
\(Q(A)=S^{-1}A\)，并把 \(A\) 在 \(Q(A)\) 中的整闭包称为其正规化。对于非约化环，本书先取 \(A_{\mathrm{red}}=A/\sqrt0\)，再正规化 \(A_{\mathrm{red}}\)，不将幂零元保留在正规化对象中。
\end{definition}
\symidx{\(Q(A)\)：约化环的全分式环}
\symidx{\(\widetilde A\)：环 \(A\) 的正规化}

非零因子乘法封闭，而且 \(A\to Q(A)\) 单射：\(a/1=0\) 意味着某个非零因子 \(s\) 满足 \(sa=0\)，从而 \(a=0\)。当 \(A\) 是整环时，全分式环就是普通分式域，因此两种定义一致。若 \(A\) 约化，则 \(A\subseteq\widetilde A\) 是整扩张。对一般 \(A\)，自然映射为
\[
A\longrightarrow A_{\mathrm{red}}\longrightarrow
\widetilde{A_{\mathrm{red}}}.
\]
第二个箭头单射，所以复合映射的核为 \(\sqrt0\)。目标环中的每个元素都在 \(A_{\mathrm{red}}\) 上整，也就是在 \(A\) 的像上整；因此这是整的环同态，但当 \(A\) 非约化时不是环包含。无论哪种情形，定义本身都不保证正规化是有限生成的 \(A\) 模。

\ProseParagraphBreak
例如，令 \(C=k[x,y]/(xy)\)。把一个商类送到它在两条坐标轴上的限制，得到单射
\[
C\longrightarrow k[x]\times k[y],\qquad
[F]\longmapsto\bigl(F(x,0),F(0,y)\bigr).
\]
每个商类唯一写成 \(c+xp(x)+yq(y)\)，所以其像恰为常数项相等的多项式对，这也说明 \(C\) 约化。这样的对 \((f,g)\) 是非零因子，当且仅当 \(f,g\) 都非零：两者非零时可在两个整环中分别约去；若 \(f=0\)，则乘以非零元素 \((x,0)\) 得到零，若 \(g=0\)，则乘以 \((0,y)\) 得到零。因此上述映射延拓为单射 \(Q(C)\to k(x)\times k(y)\)。对任意 \(F/G\in k(x)\) 和 \(P/H\in k(y)\)，其中 \(G,H\ne0\)，有
\[
\frac{(xF,0)}{(xG,y)}\longmapsto(F/G,0),\qquad
\frac{(0,yP)}{(x,yH)}\longmapsto(0,P/H).
\]
这里分子、分母都是常数项为零的多项式对，且分母的两个坐标都非零。把这两类分式相加便得到任意有理函数对，故 \(Q(C)\cong k(x)\times k(y)\)。

在这个全分式环中，\(e=(1,0)\) 满足首一方程 \(e^2-e=0\)。而且 \(C[e]=k[x]\times k[y]\)：任意 \((f,g)\) 都等于
\(e(f,f(0))+(1-e)(g(0),g)\)，右侧两对多项式属于 \(C\)。所以 \(k[x]\times k[y]\) 在 \(C\) 上整。若一个有理函数对在 \(k[x]\times k[y]\) 上整，投影到每个坐标便得到该坐标在相应多项式环上的首一方程；\(k[x],k[y]\) 都整闭，故两个坐标均为多项式。任何在 \(C\) 上整的有理函数对也在这个上环上整，因此属于它。这证明 \(k[x]\times k[y]\) 恰为 \(C\) 的正规化。原环要求两个分支上的函数在原点取相同值，正规化后这项限制消失。

对一般约化 Noether 环，极小素理想只有有限多个，并有
\[
Q(A)\cong\prod_{\mathfrak p\in\operatorname{Min}(A)}
\operatorname{Frac}(A/\mathfrak p).
\]
这一分支描述将在\cref{b3:prop:reduced-total-quotient-product}证明。对于整环，全分式环是域，正规化是其中的整环。

\subsection{非有限正规化的例子}\label{b3:subsec:0604:02}

\begin{example}\label{b3:exm:nonfinite-normalization}
取无限次代数扩域 \(K/k\)，例如 \(\overline{\mathbb Q}/\mathbb Q\)。允许正次数项的系数来自 \(K\)，却把常数项限制在 \(k\)，便得到
\[
A=k+xK[x]\subseteq B=K[x].
\]
则 \(\operatorname{Frac}(A)=K(x)\)，正规化为 \(B\)，但 \(B\) 不是有限生成的 \(A\) 模；此处 \(A\) 不是 Noether 环。
\end{example}
\begin{proof}
对每个 \(c\in K\)，有 \(c=(cx)/x\in\operatorname{Frac}(A)\)，故分式域为 \(K(x)\)。\(K\) 中每个元素在 \(k\) 上代数，因而在 \(A\) 上整；而 \(B\) 由 \(A\) 与这些常数生成，所以\cref{b3:prop:integral-closure-ring}说明 \(B/A\) 整。又 \(B=K[x]\) 是唯一分解整环，按\cref{b3:prop:ufd-integrally-closed}在 \(K(x)\) 中整闭。任何在 \(A\) 上整的 \(z\in K(x)\) 也在 \(B\) 上整，所以属于 \(B\)。这就证明了正规化等于 \(B\)。

若 \(B\) 是有限生成的 \(A\) 模，模掉同时为 \(A,B\) 理想的 \(xK[x]\)，就会得到 \(K=B/xK[x]\) 是有限维 \(A/xK[x]=k\) 向量空间，与 \([K:k]=\infty\) 矛盾。最后，理想 \(\mathfrak m=xK[x]\) 满足
\(\mathfrak m/\mathfrak m^2\cong K\)，其上 \(A\) 的作用经 \(k\) 因子化；若 \(\mathfrak m\) 有限生成，此商将是有限维 \(k\) 空间，仍矛盾。因此 \(A\) 确实不是 Noether 环。
\end{proof}

这个例子说明没有额外条件时正规化未必有限；其底环是非 Noether 环。域上有限型代数拥有更强的结构；\cref{b3:thm:normalization-finite}将证明这一范围内的正规化有限，任意基域都允许。

\subsection{尖点与节点曲线的正规化}\label{b3:subsec:0604:03}

\begin{proposition}\label{b3:prop:cusp-normalization}
设 \(k\) 为任意域。映射 \(x\mapsto t^2\)、\(y\mapsto t^3\) 给出同构
\[
k[x,y]/(y^2-x^3)\cong k[t^2,t^3].
\]
此整环的分式域是 \(k(t)\)，正规化为 \(k[t]\)，并且 \(k[t]\) 由 \(1,t\) 作为 \(k[t^2,t^3]\) 模生成。
\end{proposition}
\begin{proof}
按首一多项式 \(y^2-x^3\) 除法，每个商类有代表 \(P(x)+y\cdot Q(x)\)。代入后为 \(P(t^2)+t^3\cdot Q(t^2)\)，前一项只含偶次幂，后一项只含不小于三的奇次幂，故等于零恰好要求 \(P=Q=0\)。因此核就是所列主理想，像按定义为 \(k[t^2,t^3]\)。

元素 \(t=t^3/t^2\) 在分式域中，且满足 \(T^2-t^2=0\)，所以 \(k[t]=A+At\) 在 \(A=k[t^2,t^3]\) 上有限整。\(k[t]\) 在 \(k(t)\) 中整闭，任何在 \(A\) 上整的分式也在 \(k[t]\) 上整，因而属于 \(k[t]\)。于是它正是整闭包。
\end{proof}

这里 \(1,t\) 是模生成元而不是模基，因为 \(t^2\cdot t-t^3\cdot1=0\) 是一个非平凡的 \(A\) 系数关系。两项模生成式证明了有限性，而分式域相等说明双有理性；相应的分式域扩张次数为一，并不是模生成元的个数二。

\begin{proposition}\label{b3:prop:node-normalization}
设 \(k\) 为特征不等于 \(2\) 的域。映射
\[
x\longmapsto t^2-1,\qquad y\longmapsto t(t^2-1)
\]
给出同构
\[
k[x,y]/\bigl(y^2-x^2(x+1)\bigr)
 \cong A=k\bigl[t^2-1,t(t^2-1)\bigr].
\]
其分式域为 \(k(t)\)，正规化为 \(B=k[t]\)，且 \(B=A+At\)。这称为本章的\term[jiedian]{节点}[node]模型。
\end{proposition}
\begin{proof}
参数值 \(t=1\) 和 \(t=-1\) 都给出 \(x=y=0\)，所以这两个不同参数值落在同一个原点。只要 \(x\ne0\)，便可由 \(t=y/x\) 恢复参数。

模掉关于 \(y\) 的首一关系后，每个商类写为 \(P(x)+y\cdot Q(x)\)。像为
\[
P(t^2-1)+t(t^2-1)\cdot Q(t^2-1).
\]
第一项只含偶次幂，第二项只含奇次幂。两者之和为零时分别为零；由于 \(t^2-1\) 在 \(k[t]\) 中超越且非零，得到 \(P=Q=0\)。这证明核的断言。分式 \(y/x=t\) 给出分式域相等，也说明 \(A_x=B_x\)。方程 \(t^2=x+1\) 给出 \(B=A+At\)，最后由 \(k[t]\) 的整闭性判定它就是正规化。
\end{proof}

特征不为二用于保证 \(t=1\) 与 \(t=-1\) 是两个不同点；上面的核及整闭包计算本身不需除以二。这个节点的原环还有函数值描述
\[
A=\bigl\{\,f\in k[t]\mid f(1)=f(-1)\,\bigr\}.
\]
\Cref{b3:prop:node-conductor}将证明这一等式，并由此计算正规化商模和导子理想。

\subsection{整闭包、有限性与双有理性}\label{b3:subsec:0604:04}

设 \(A\subseteq B\) 为交换含幺整环的包含。若诱导的分式域映射 \(\operatorname{Frac}(A)\to\operatorname{Frac}(B)\) 是同构，则称 \(B/A\) 为\term[shuangyouli-kuozhang]{双有理扩张}[birational extension]；相应的仿射映射称为双有理。正规化按定义双有理且整，但有限性是额外结论。例如前面的非 Noether 例仍双有理而不有限。

\ProseParagraphBreak
反过来，有限整扩张不必双有理：\(k[t^2]\subseteq k[t]\) 虽然由 \(1,t\) 生成，却诱导二次分式域扩张 \(k(t)/k(t^2)\)。即使特征为二、底层点集可能不能区分平方根，这个域扩张仍有次数二。双有理也不意味着有限，例如 \(k[t]\subseteq k[t,t^{-1}]\) 分式域相同，却把零点处的分母倒置；\(t^{-1}\) 不在整闭环 \(k[t]\) 上整，所以此扩张不有限。

\ProseParagraphBreak
尖点与节点的有限性在\cref{b3:prop:cusp-normalization}及\cref{b3:prop:node-normalization}中都体现为两项模生成式。一般有限型代数未必有这样显式的参数表达；下一节借 Noether 正规化和有限域扩张的结构证明其正规化仍然有限。

% !TeX root = ../../Book3.tex
% 纲要标题：### 6.5 正规化有限性的证明
% 纲要位置：工作记录/计划纲要.md，第 1989 行。
% 纲要细目（写作范围）：
% * 用 Noether 正规化把有限型整环的整闭包问题归约为多项式环在有限分式域扩张中的整闭包问题
% * 对有限可分扩张建立迹配对和公共分母控制引理，利用多项式环的整闭性与 Noether 性得到有限生成
% * 正特征下单独建立不可分情形所需的有限性引理；通过有限个系数和纯不可分扩张的代数处理完成证明，不把可分迹论证无条件套用，也不默设基域完美
% * 从整环推广到约化有限型代数的有限个极小素分支；明确全分式环中的正规化与各分支整闭包的关系
% * 本节接续第6.4节的正规化对象，准确陈述并完整证明有限性定理；不以附录中的一般 Nagata 环定理代替所需引理
% * Nagata 性、卓越性与正规化有限性的适用范围

\section{正规化有限性的证明}
\label{b3:sec:0605}

上一节定义了正规化，并在曲线中找到了有限个新生成元。一般有限型代数是否也只需添有限个元素，不能由定义直接得出。本节的关键是控制整元素的分母：可分扩张中用迹配对，不可分扩张中先进行一个只涉及有限个系数的纯不可分变基。整个论证允许基域不完美。这一有限性论证可对照 \cite[\S33, Integral closure of a Noetherian integral domain, pp. 261--263]{Matsumura1989CommutativeRingTheory}；下文把其中的不可分变基与分支处理分别展开。

\begin{figure}[htbp]
\centering
\begin{tikzcd}[column sep=6em,row sep=2.5em]
\text{有限型 }k\text{ 整环}
  \arrow[r,"\text{Noether 正规化}"]
  & A=k[x_1,\ldots,x_d]
    \arrow[dl,"\text{有限可分扩张}"']
    \arrow[d,"\text{有限不可分扩张}"] \\
\text{迹配对与公共分母}\arrow[dr]
  & \text{有限纯不可分变基}
    \arrow[l,"\text{化为可分}"'] \\
  & \text{整环情形的整闭包有限}\arrow[d] \\
  & \text{有限型约化环逐极小素分支处理}
\end{tikzcd}
\caption[正规化有限性的证明结构]{正规化有限性的证明结构。基域 \(k\) 任意，记 \(K=\operatorname{Frac}(A)\)，两条分支分别处理有限扩张 \(L/K\) 的可分与不可分情形。}
\label{b3:fig:normalization-finiteness-proof}
\end{figure}

\cref{b3:fig:normalization-finiteness-proof}中，迹配对始终用于有限可分扩张；不可分情形先作有限变基，再通过有限模的传递及 Noether 子模性质回到原环。约化环则需要全分式环的极小素分支分解。\secref{b3:sec:0704}的交叉直线计算及\chapref{b3:ch:20}的 Serre 正规性判据也用到这一分解。\footnote{塞尔（Jean-Pierre Serre，1926-），法国数学家，研究代数拓扑、代数几何与数论。}

\subsection{由 Noether 正规化归约整闭包问题}
\label{b3:subsec:0605:01}

设 \(C\) 为域 \(k\) 上的有限型整环。由\cref{b3:thm:noether-normalization}，存在代数独立的 \(x_1,\ldots,x_d\in C\)，使
\[
A=k[x_1,\ldots,x_d]\subseteq C
\]
为有限扩张。记 \(K=\operatorname{Frac}(A)\)、\(L=\operatorname{Frac}(C)\)。\(C\) 的有限组代数生成元在 \(K\) 上代数，且生成 \(L\)，因此 \([L:K]<\infty\)。

\ProseParagraphBreak
更一般地，若 \(L'/L\) 是任意有限域扩张，则 \(C\) 在 \(L'\) 中的整闭包等于 \(A\) 在 \(L'\) 中的整闭包。一个方向由 \(A\subseteq C\) 直接得到；另一个方向用 \(C/A\) 整和\cref{b3:prop:integral-transitivity}。所以只须证明：多项式环 \(A\) 在任意有限分式域扩张中的整闭包是有限生成的 \(A\) 模。若这一步完成，有限组 \(A\) 模生成元也生成同一个 \(C\) 模，因为 \(A\subseteq C\)。

\subsection{可分扩张中的迹配对与公共分母}
\label{b3:subsec:0605:02}

\begin{lemma}\label{b3:lem:normalization-separable-finite}
设 \(A\) 为交换含幺整闭 Noether 整环，\(K=\operatorname{Frac}(A)\)，\(L/K\) 为有限可分扩张。则 \(A\) 在 \(L\) 中的整闭包 \(B\) 是有限生成的 \(A\) 模。
\end{lemma}
\begin{proof}
先取一组 \(K\) 基 \(v_1,\ldots,v_n\)；乘以适当的非零 \(A\) 元素，可使各基元素对 \(A\) 整。确切地，若
\[
v_i^m+c_1v_i^{m-1}+\cdots+c_m=0,\qquad c_j\in K,
\]
选非零 \(d_i\in A\) 清除全部 \(c_j\) 的分母，则 \(u_i=d_iv_i\) 满足首一方程
\[
u_i^m+d_ic_1u_i^{m-1}+\cdots+d_i^mc_m=0
\]
且其系数均在 \(A\) 内。缩放后 \(u_1,\ldots,u_n\) 仍为 \(K\) 基。

对任一整元素 \(z\in L\)，域迹 \(\operatorname{Tr}_{L/K}(z)\) 属于 \(A\)。事实上，第一卷定理18.2.3把可分扩张的迹写成各 \(K\) 嵌入像之和；每个像都满足 \(z\) 的同一个首一 \(A\) 系数方程，故都对 \(A\) 整。它们的和也整，而迹又属于 \(K\)，由 \(A\) 整闭得此结论。

令
\[
T=\bigl(\operatorname{Tr}_{L/K}(u_i u_j)\bigr)_{1\le i,j\le n},
\qquad \Delta=\det T.
\]
各 \(u_i u_j\) 都整，故 \(T\) 的条目在 \(A\) 中。第一卷定理18.3.3说明有限可分扩张的迹配对非退化，所以 \(\Delta\ne0\)。给定 \(b\in B\)，写 \(b=\sum_i\lambda_i u_i\)、\(\lambda_i\in K\)。各 \(bu_j\) 仍整，于是
\[
T(\lambda_1,\ldots,\lambda_n)^{\mathsf T}
=\bigl(\operatorname{Tr}_{L/K}(bu_1),\ldots,
        \operatorname{Tr}_{L/K}(bu_n)\bigr)^{\mathsf T}\in A^n.
\]
用伴随矩阵乘上式，得每个 \(\Delta\lambda_i\in A\)。因此
\begin{equation}\label{b3:eq:normalization-trace-denominator}
B\subseteq\Delta^{-1}\sum_{i=1}^nAu_i.
\end{equation}
因为 \(u_1,\ldots,u_n\) 是 \(K\) 基，它们也在 \(A\) 上线性无关，故 \(\sum_iAu_i\cong A^n\)。乘以非零标量 \(\Delta^{-1}\) 不改变这个模的同构类型，所以右端是有限自由 \(A\) 模，左端由\cref{b3:prop:integral-closure-ring}是一个 \(A\) 子模。因 \(A\) 为 Noether 环，\cref{b3:prop:noether-finite-modules}保证其子模 \(B\) 仍有限生成。
\end{proof}

这一步真正使用可分性的地方是 \(\Delta\ne0\)。在非平凡纯不可分扩张中，域迹恒为零，所有这样的迹矩阵都退化；不能在此情形下仍写 \(\Delta^{-1}\)。下一小节通过有限变基把问题移到可以使用这个引理的地方。

\subsection{正特征不可分扩张中的有限性引理}
\label{b3:subsec:0605:03}

\begin{lemma}\label{b3:lem:finite-purely-inseparable-base-change}
设 \(k\) 为特征 \(p>0\) 的域，\(d\ge0\)，\(x_1,\ldots,x_d\) 为代数独立元，\(A=k[x_1,\ldots,x_d]\)，\(K=\operatorname{Frac}(A)\)，\(L/K\) 为有限扩张。在同一个代数闭包内，可以选取整数 \(e\ge0\)、\(q=p^e\)、有限纯不可分扩张 \(k'/k\) 及元素 \(y_j\) 满足 \(y_j^q=x_j\)，使
\[
K'=k'(y_1,\ldots,y_d),\qquad L'=LK'
\]
满足 \(L'/K'\) 有限可分。并且 \(A'=k'[y_1,\ldots,y_d]\) 是多项式环，是 \(A=k[x_1,\ldots,x_d]\) 上的有限生成模。
\end{lemma}
\begin{proof}
取有限组域生成元 \(z_1,\ldots,z_s\)，使 \(L=K(z_1,\ldots,z_s)\)。将每个首一不可约最小多项式写成
\[
f_i(T)=g_i(T^{q_i}),\qquad q_i=p^{e_i},\qquad g_i'(T)\ne0.
\]
这可以通过反复除去全部指数的公因子 \(p\) 实现。\(g_i\) 仍不可约，否则代入 \(T^{q_i}\) 会分解 \(f_i\)；因而 \(g_i\) 可分。

取 \(q=p^e\) 为所有 \(q_i\) 的公倍数。将各 \(g_i\) 的有限多个系数写成 \(x_1,\ldots,x_d\) 的有理函数，只涉及有限多个 \(k\) 中的分子、分母多项式系数，设其集合为 \(E\)。令
\[
k'=k\bigl(\,c^{1/q}:c\in E\,\bigr),\qquad y_j=x_j^{1/q}.
\]
各次根在固定代数闭包中唯一，\(k'/k\) 有限纯不可分，且 \((k')^q\subseteq k\)。因此 \(y_1,\ldots,y_d\) 在 \(k'\) 上代数独立：若有非零多项式关系 \(\sum_\alpha b_\alpha y^\alpha=0\)，将其取 \(q\) 次幂，Frobenius 加法性给出 \(\sum_\alpha b_\alpha^q x^\alpha=0\)。这里每个 \(b_\alpha^q\in k\)，而 Frobenius 在域 \(k'\) 上单射，至少一个非零系数的 \(q\) 次幂仍非零；这便是 \(x_1,\ldots,x_d\) 在 \(k\) 上的非零关系，矛盾。

每个 \(g_i\) 的系数 \(c\in K\) 都有 \(q_i\) 次根属于 \(K'\)。确切地，写 \(c=P(x)/Q(x)\)，其中 \(Q\ne0\)，并将 \(P\) 按多重指数写成 \(P(x)=\sum_\alpha a_\alpha x^\alpha\)。各 \(a_\alpha\) 的 \(q\) 次根已添入 \(k'\)，所以
\[
a_\alpha^{1/q_i}
  =\bigl(a_\alpha^{1/q}\bigr)^{q/q_i}\in k',
\qquad
P_i^\sharp(y)
  =\sum_\alpha a_\alpha^{1/q_i}y^{(q/q_i)\alpha}.
\]
Frobenius 加法性给出 \((P_i^\sharp(y))^{q_i}=P(x)\)。对 \(Q\) 同样构造 \(Q_i^\sharp(y)\)，则 \((Q_i^\sharp(y))^{q_i}=Q(x)\ne0\)，所以它可以作分母，且
\[
c^{1/q_i}=\dfrac{P_i^\sharp(y)}{Q_i^\sharp(y)}\in K'.
\]
写 \(g_i(T)=\sum_jc_{ij}T^j\)，令
\[
h_i(T)=\sum_jc_{ij}^{1/q_i}T^j\in K'[T].
\]
则 \(h_i(T)^{q_i}=f_i(T)\)，故 \(h_i(z_i)=0\)。而 \(h_i\) 可分：若 \(w\) 同时零化 \(h_i\) 与 \(h_i'\)，取 \(q_i\) 次幂后即得
\[
g_i(w^{q_i})=0,\qquad g_i'(w^{q_i})=0,
\]
其中导数系数中的整数属于素域，取 \(q_i\) 次幂后不变。这与 \(g_i\) 可分矛盾。\(z_i\) 在 \(K'\) 上的最小多项式整除 \(h_i\)，而可分多项式的不可约因子仍可分，故各 \(z_i\) 在 \(K'\) 上可分。它们共同生成的 \(L'/K'\) 也是有限可分扩张。

最后取 \(k'/k\) 的有限基 \(c_1,\ldots,c_t\)。因为 \(y_j^q=x_j\)，\(A'\) 由有限组元素
\[
c_l y_1^{a_1}\cdots y_d^{a_d},
\qquad 1\le l\le t,\quad 0\le a_j<q,
\]
作为 \(A\) 模生成。\(d=0\) 或所有 \(q_i=1\) 时，上述构造分别采用空变量组或 \(q=1\)，同样成立。
\end{proof}

只取“实际用到的有限多个系数”的次根，是这个引理中的有限性所在。任意不完美域的扩张 \(k^{1/p}/k\) 可能无限，因而直接把全部 \(p\) 次根都添进去，反而会失去这里需要的有限生成模控制。

\begin{theorem}\label{b3:thm:polynomial-normalization-finite}
设 \(k\) 为任意域，\(d\ge0\)，\(A=k[x_1,\ldots,x_d]\)，\(K=\operatorname{Frac}(A)\)。对任意有限扩张 \(L/K\)，\(A\) 在 \(L\) 中的整闭包是有限生成的 \(A\) 模。
\end{theorem}
\begin{proof}
多项式环 \(A\) 是 Noether 唯一分解整环，因而由\cref{b3:prop:ufd-integrally-closed}整闭。若 \(L/K\) 可分，直接应用\cref{b3:lem:normalization-separable-finite}。特征零的情形已全部包括在内。

在正特征中，按\cref{b3:lem:finite-purely-inseparable-base-change}构造 \(A'\)、\(K'\)、\(L'\)。\(A'\) 也是整闭 Noether 多项式环，而 \(L'/K'\) 可分，所以 \(A'\) 在 \(L'\) 中的整闭包 \(B'\) 有限于 \(A'\)。由于 \(A'\) 又有限于 \(A\)，\(B'\) 是有限生成的 \(A\) 模。原来每个对 \(A\) 整的 \(b\in L\) 也对 \(A'\) 整，因此原整闭包 \(B\) 包含于 \(B'\)。\(B\) 是 \(A\) 子模，\(A\) 的 Noether 性于是保证 \(B\) 有限。
\end{proof}

\subsection{约化代数的极小素分支与全分式环}
\label{b3:subsec:0605:04}

整环只有一个分式域；约化环可能有多个极小素分支，故应在全分式环中同时处理各分支。下面先给出此前使用的乘积描述的证明。

\begin{proposition}\label{b3:prop:reduced-total-quotient-product}
设 \(C\ne0\) 为交换含幺约化 Noether 环，极小素理想为 \(\mathfrak p_1,\ldots,\mathfrak p_r\)，\(K_i=\operatorname{Frac}(C/\mathfrak p_i)\)。则自然映射给出同构
\begin{equation}\label{b3:eq:reduced-total-quotient-product}
Q(C)\cong\prod_{i=1}^rK_i.
\end{equation}
若 \(B_i\) 为 \(C/\mathfrak p_i\) 在 \(K_i\) 中的整闭包，则 \(C\) 在 \(Q(C)\) 中的整闭包为 \(\prod_iB_i\)。
\end{proposition}
\begin{proof}
因为 \(C\) 同时约化且 Noether，\cref{b3:prop:reduced-ring-no-embedded-primes}给出 \(\operatorname{Ass}_C(C)=\min\operatorname{Spec}C\)。再由\cref{b3:thm:associated-zero-divisors}，\(C\) 的非零因子组成
\[
S=C\setminus\bigcup_i\mathfrak p_i.
\]
约化性给出 \(\bigcap_i\mathfrak p_i=0\)，故 \(C\rightarrow\prod_iK_i\) 单射。每个 \(s\in S\) 在所有 \(K_i\) 中非零，因此该映射延伸为单射 \(Q(C)=S^{-1}C\rightarrow\prod_iK_i\)。

为证明满射，对每个 \(i\) 选
\[
c_i\in\bigcap_{j\ne i}\mathfrak p_j\setminus\mathfrak p_i.
\]
这样的 \(c_i\) 可由逐个取 \(a_j\in\mathfrak p_j\setminus\mathfrak p_i\) 后相乘得到；\(r=1\) 时取 \(c_1=1\)。给定第 \(i\) 个域中的分式 \(\bar a/\bar b\)，其中 \(a,b\in C\)、\(b\notin\mathfrak p_i\)，令
\[
s=bc_i+\sum_{j\ne i}c_j.
\]
它模 \(\mathfrak p_i\) 等于非零的 \(bc_i\)，模每个 \(\mathfrak p_j\)、\(j\ne i\) 等于非零的 \(c_j\)，所以 \(s\in S\)。分式 \(ac_i/s\) 的第 \(i\) 个坐标就是 \(\bar a/\bar b\)，其余坐标全零。逐坐标相加即可得到任意元组，满射得证。

若 \(b\in Q(C)\) 对 \(C\) 整，投影它的首一方程可见每个坐标 \(b_i\) 属于 \(B_i\)。反之，若各 \(b_i\in B_i\)，为每个 \(i\) 取一个零化 \(b_i\) 的首一 \((C/\mathfrak p_i)\) 系数多项式，并把系数提升为首一 \(f_i(T)\in C[T]\)。对元组 \(b=(b_i)_i\)，多项式 \(\prod_i f_i(T)\) 在每个坐标上都有一个为零的因子 \(f_i(b_i)=0\)，故 \(\prod_i f_i(b)=0\)。这个首一方程证明 \(b\) 对 \(C\) 整。
\end{proof}

特别地，全分式环中的不同分支彼此独立，正规化也逐分支进行。这里的有限个分支来自 Noether 性；约化性则保证环嵌入这些分式域的乘积。若原环非约化，本书按\cref{b3:def:normalization}的约定先取约化商，再谈这一正规化。

\subsection{正规化有限性定理的完整证明}
\label{b3:subsec:0605:05}

\begin{theorem}\label{b3:thm:normalization-finite}
设 \(k\) 为任意域。
\begin{enumerate}
\item 若 \(C\) 为交换含幺有限型 \(k\) 整环，\(L'/\operatorname{Frac}(C)\) 为有限域扩张，则 \(C\) 在 \(L'\) 中的整闭包是有限生成的 \(C\) 模。特别地，\(C\) 的正规化有限。
\item 若 \(C\) 为交换含幺约化有限型 \(k\) 代数，则它在全分式环中的正规化是有限生成的 \(C\) 模。
\end{enumerate}
\end{theorem}
\begin{proof}
第一项取\cref{b3:thm:noether-normalization}给出的有限扩张 \(A=k[x_1,\ldots,x_d]\subseteq C\)。\subsecref{b3:subsec:0605:01}已经核验：\(L'/\operatorname{Frac}(A)\) 有限，且 \(A\) 与 \(C\) 在 \(L'\) 中的整闭包相同。\cref{b3:thm:polynomial-normalization-finite}保证此整闭包有限于 \(A\)，故也有限于 \(C\)。

第二项中，\(C=0\) 时结论直接成立。否则 \(C\) 是 Noether 环，有有限个极小素理想 \(\mathfrak p_i\)。每个 \(C/\mathfrak p_i\) 都是有限型 \(k\) 整环，由第一项其正规化 \(B_i\) 为有限生成的 \(C/\mathfrak p_i\) 模，从而也是有限生成的 \(C\) 模。\cref{b3:prop:reduced-total-quotient-product}把 \(C\) 的正规化识别为有限乘积 \(\prod_iB_i\)，所以它仍是有限生成的 \(C\) 模。
\end{proof}

可分性只是证明中的工具，最终定理允许任意基域及有限分式域扩张。上述有限性也为下一节导子理想中的有限环包含提供依据。

\subsection{Nagata 性、卓越性与正规化的有限性}\label{b3:subsec:0605:06}
\label{b3:subsec:0604:05}

上述证明先利用有限型性选出多项式子环 \(A\)，使原整环在 \(A\) 上有限；分式域扩张的有限性再保证只需处理有限组域生成元及其最小多项式中的有限个系数。一般 Noether 整环未必有这样的多项式子环，其正规化也未必有限。

\ProseParagraphBreak
后续资料中常见的\term[nagata-huan]{Nagata 环}[Nagata ring]，\footnote{永田雅宜（假名：ながた まさよし，罗马音：Masayoshi Nagata，1927-2008），日本数学家，研究交换代数与代数几何，构造 Hilbert 第十四问题的反例。}按一种标准定义，是满足下列条件的 Noether 环 \(A\)：对每个素理想 \(\mathfrak p\)，以及 \(\operatorname{Frac}(A/\mathfrak p)\) 的每个有限域扩张 \(L\)，\(A/\mathfrak p\) 在 \(L\) 中的整闭包都是有限生成的 \(A/\mathfrak p\) 模。\cref{b3:thm:normalization-finite}对这些商整环及有限域扩张均适用，因而域上有限型代数都是 Nagata 环。\term[zhuoyue-huan]{卓越环}[excellent ring]还包含关于形式纤维、维数链及正则性的条件，其定义及基本性质见\secref{b3:sec:D04}。

% !TeX root = ../../Book3.tex
% 纲要标题：### 6.6 导子理想与曲线正规化的比较
% 纲要位置：工作记录/计划纲要.md，第 1997 行。
% 纲要细目（写作范围）：
% * 对有限双有理扩张 \(A\subseteq B\) 定义导子理想 \(\mathfrak c=\{a\in A:aB\subseteq A\}\)，说明其同时为 \(A\) 和 \(B\) 的理想
% * 在 \(A=k[t^2,t^3]\subseteq B=k[t]\) 中计算 \(\mathfrak c=t^2k[t]\)、\(B/A\) 与奇点处的差异
% * 以特征不为2的曲线 \(y^2=x^2(x+1)\) 为节点模型，使用 \(x=t^2-1\)、\(y=t(t^2-1)\) 核验正规化及奇点上的两个原像
% * 比较有限双有理映射、点集上的行为和局部环的改变；平坦性与非分歧性的独立检验留到第23.4节
% ---

\section{导子理想与曲线正规化的比较}\label{b3:sec:0606}

正规化中哪些元素原来就能相乘而不离开旧环？这个问题由导子理想回答。它既是旧环的理想，也是新环的理想，因此可以直接比较两侧的局部化与商模。本节继续使用尖点和节点的显式环，不以一般奇点理论解释尚未计算的现象。

\subsection{有限双有理扩张中的导子理想}\label{b3:subsec:0606:01}

\begin{definition}\label{b3:def:conductor}
设 \(A\subseteq B\) 为交换含幺环的包含。定义其\term[daozi-lixiang]{导子理想}[conductor ideal]
\[
\mathfrak c_{B/A}
 =\{\,a\in A\mid aB\subseteq A\,\}
 =\operatorname{Ann}_A(B/A).
\]
本节主要考虑整环的有限双有理扩张。
\end{definition}
\symidx{\(\mathfrak c_{B/A}\)：扩张 \(A\subseteq B\) 的导子理想}

\begin{proposition}\label{b3:prop:conductor-basic}
设 \(A\subseteq B\) 为交换含幺环的包含，\(\mathfrak c_{B/A}=\{a\in A:aB\subseteq A\}\) 为导子理想。则 \(\mathfrak c_{B/A}\) 是包含在 \(A\) 中的最大 \(B\) 理想，因而同时为 \(A,B\) 的理想。若整环扩张 \(B/A\) 有限且双有理，则存在非零 \(d\in A\) 使 \(dB\subseteq A\)。对每个 \(c\in\mathfrak c_{B/A}\)，都有 \(A_c=B_c\)。
\end{proposition}
\begin{proof}
若 \(a\in\mathfrak c_{B/A}\)、\(b\in B\)，则 \(ab\in A\)，且 \((ab)B\subseteq aB\subseteq A\)，所以 \(ab\) 也在导子理想中。任意 \(B\) 理想 \(J\subseteq A\) 满足 \(jB\subseteq J\subseteq A\)，故 \(J\subseteq\mathfrak c_{B/A}\)。

有限双有理情形选 \(B\) 的有限组 \(A\) 模生成元 \(b_i=u_i/v_i\)，其中 \(u_i,v_i\in A\)、\(v_i\ne0\)。令 \(d=\prod_i v_i\)，便有 \(db_i\in A\)，于是 \(dB\subseteq A\) 且 \(d\ne0\)。最后对 \(b\in B\)，有 \(b=(cb)/c\in A_c\)，这证明 \(B_c\subseteq A_c\)，反向包含由原包含得到。
\end{proof}

对有限扩张，差异还可以用支集精确表述：
\[
\operatorname{Supp}_A(B/A)=V_A(\mathfrak c_{B/A}).
\]
这里 \(B/A\) 是有限生成模，故\cref{b3:thm:support-finite}把其支集识别为零化子的零点集。因而一个素理想 \(\mathfrak p\) 不含导子理想，当且仅当 \(A_{\mathfrak p}\to B\otimes_AA_{\mathfrak p}\) 是同构；这只是局部环比较，不是光滑性判据。

\subsection{尖点的导子理想与正规化商模}\label{b3:subsec:0606:02}

\begin{proposition}\label{b3:prop:cusp-conductor}
设 \(k\) 为域，\(A=k[t^2,t^3]\subseteq B=k[t]\)。则
\[
A=k\oplus t^2k[t],\qquad
B/A\cong k\cdot[t],\qquad
\mathfrak c_{B/A}=t^2k[t]=(t^2,t^3)_A.
\]
其中 \(B/A\) 上 \(A\) 的作用经 \(A/(t^2,t^3)\cong k\) 因子化。
\end{proposition}
\begin{proof}
所有不小于二的整数都由二和三作非负整数线性组合，故 \(A\) 恰由没有一次项的多项式组成。于是商空间由 \([t]\) 生成，且 \(t^2,t^3\) 都零化它。若 \(a=c+t^2\cdot h(t)\in A\)，则
\(at=ct+t^3\cdot h(t)\) 属于 \(A\) 当且仅当 \(c=0\)。由于 \(B=A+At\)，条件 \(aB\subseteq A\) 等价于 \(at\in A\)，故导子为 \(t^2k[t]\)。最后任意次数至少四的单项式可除以 \(t^2\) 而仍属于 \(A\)，再加上 \(t^2,t^3\) 两项，得到其作为 \(A\) 理想的生成式。
\end{proof}

导子只在尖点极大理想 \(\mathfrak m=(t^2,t^3)\) 处消失，离开该点两环局部化相同。该点上只有一个原像 \((t)\subseteq k[t]\)，但纤维代数是
\[
B/\mathfrak mB=k[t]/(t^2),
\]
并不是 \(k\)。所以“只有一个原像点”仍会遗漏非约化纤维。局部化后的差商
\(B_{(t)}/A_{\mathfrak m}\) 仍为一维 \(k\) 空间；这里两个局部化嵌在 \(k(t)\) 中，并且 \(B\otimes_AA_{\mathfrak m}=B_{(t)}\)。后一个等式可由\cref{b3:lem:integral-field-criterion}与局部化素理想对应看出：左环的极大理想收缩为 \(A_{\mathfrak m}\) 的唯一极大理想，因此左环也只有极大理想 \((t)\)，因而所有不属于 \((t)\) 的多项式都成为单位。

\subsection{节点正规化与奇点上的两个原像}\label{b3:subsec:0606:03}

\begin{proposition}\label{b3:prop:node-conductor}
设 \(k\) 为特征不等于 \(2\) 的域，令
\[
A=k\bigl[t^2-1,t(t^2-1)\bigr]\subseteq B=k[t].
\]
则
\[
A=k+(t^2-1)k[t]
 =\bigl\{\,f\in k[t]\mid f(1)=f(-1)\,\bigr\},
\]
并有
\[
B/A\cong k,\qquad
\mathfrak c_{B/A}=(t^2-1)k[t]=(x,y)_A,
\]
其中 \(x=t^2-1\)、\(y=t(t^2-1)\)。原点的两个原像为 \((t-1)\) 和 \((t+1)\)。
\end{proposition}
\begin{proof}
记 \(h=t^2-1\)。若 \(j\) 为偶数，则 \(ht^j=h(h+1)^{j/2}\in k[h]\)；若 \(j\) 为奇数，则
\(ht^j=ht(h+1)^{(j-1)/2}\in k[h,ht]\)。
故 \(hB\subseteq A\)，而 \(A\) 的生成元都属于 \(k+hB\)，得到第一项等号。

任意 \(f\in k[t]\) 除以 \(h\) 的余式为 \(a+bt\)。条件 \(f(1)=f(-1)\) 等价于 \(2b=0\)，由特征假设等价于 \(b=0\)，这证明第二项等号。线性映射 \(f\mapsto f(1)-f(-1)\) 满射到 \(k\)，核为 \(A\)，所以 \(B/A\) 一维；\(A\) 的作用经公共值 \(a(1)=a(-1)\) 给出。

因为 \(hB\subseteq A\)，导子包含 \(hB\)。若 \(a\) 在导子中，则 \(a\in A\) 且 \(at\in A\)，从而
\(a(1)=a(-1)\) 与 \(a(1)=-a(-1)\) 同时成立。因二可逆，二者均为零，所以 \(h\mid a\)。这证明导子等于 \(hB\)。前面对 \(ht^j\) 的表达又说明其作为 \(A\) 理想由 \(h,ht\) 生成，即 \((x,y)\)。

最后 \((x,y)B=hB\)，而 \(h=(t-1)(t+1)\) 有两个不同线性因子。商环及中国剩余定理给出
\[
B/(x,y)B\cong k[t]/(t^2-1)\cong k\times k.
\]
因此其两个素理想正给出所述两个原像。
\end{proof}

\subsection{有限双有理映射的点集与局部环}\label{b3:subsec:0606:04}

尖点和节点的正规化都有限且双有理，差商的 \(k\) 维数也都为一；它们仍以不同方式改变特殊点。尖点只增加同一点附近缺失的函数，节点则把同一个原点上的两种参数值分开。前者纤维为 \(k[t]/(t^2)\)，后者纤维为 \(k\times k\)，所以点的个数与纤维的幂零结构必须同时保留。

\ProseParagraphBreak
对节点令 \(\mathfrak m=(x,y)\)、\(S=A\setminus\mathfrak m\)。局部正规化
\[
B\otimes_AA_{\mathfrak m}=S^{-1}k[t]
\]
有两个极大理想，分别来自 \((t-1)\) 与 \((t+1)\)，是一个半局部整环。它并不是两个局部环的直积：左侧仍嵌在域 \(k(t)\) 中，而非零环的二项直积有零因子。再分别在这两个极大理想处局部化，才得到
\[
k[t]_{(t-1)},\qquad k[t]_{(t+1)}.
\]
中国剩余分解出现在模掉 \(\mathfrak m\) 的纤维中，不能提前搬到整个局部正规化环上。

\ProseParagraphBreak
在两例中，导子理想之外的局部化都是同构；在特殊点，有限性与双有理性则不足以决定平坦性和非分歧性。它们需要对模或微分再作独立检验，留到\secref{b3:sec:2304}进行。本章没有把正规化称为这些更强性质的实例。


\begin{chaptersummary}
有限型 \(k\) 代数的极大理想 \(\mathfrak m\) 具有有限次剩余域扩张，即 \([A/\mathfrak m:k]<\infty\)。当 \(k\) 代数闭时，结构映射 \(k\to A/\mathfrak m\) 就是同构，于是闭点由坐标值表示；强零点定理进一步说明，普通零点集能恢复根理想，不能恢复被幂零元记录的加厚信息。

\ProseParagraphBreak
正规化是在全分式环中取整闭包，非约化对象先取约化。域上有限型代数的有限性证明先借 Noether 正规化得到多项式底环：可分部分由非退化迹配对控制公共分母，不可分部分则只添有限多个系数的根及有限个参数的根。约化代数再按有限个极小素分支合并；没有假设基域完美。

\ProseParagraphBreak
对有限双有理整环扩张 \(A\subseteq B\)，导子是包含于 \(A\) 的最大 \(B\) 理想，也是 \(A\) 模 \(B/A\) 的零化子。尖点与节点的差商都只有一维，前者在特殊点给出非约化的单点纤维，后者给出两个点。节点的局部正规化仍是半局部整环，不能把纤维的直积分解当成整个局部环的直积分解。
\end{chaptersummary}

\BookExercises

\begin{exercise}\label{b3:ex:06-rational-field-algebra}
设 \(k\) 为域，\(t\) 为变量，\(0\ne D\in k[t]\)，令 \(R=k[t,1/D]\)。把 \(D\) 写成 \(c\prod_{i=1}^s p_i^{e_i}\)，其中 \(c\in k^\times\)、\(e_i\ge1\)，\(p_i\) 是互不相同的首一不可约多项式；\(D\) 为常数时取 \(s=0\)。求 \(R\) 的全部单位，并证明对任意不可约多项式 \(h\)，有 \(1/h\in R\) 当且仅当 \(h\mid D\)。说明为何 \(R\) 是有限型 \(k\) 代数且 \(\operatorname{Frac}(R)=k(t)\)，但 \(R\ne k(t)\)。
\end{exercise}
\begin{solution}
每个 \(p_i\) 在 \(R\) 中可逆，因为 \(1/p_i=(D/p_i)/D\)。若 \(F/D^n\) 为单位，写其逆为 \(G/D^m\)，便有 \(FG=D^{n+m}\)。由 \(k[t]\) 的唯一分解，\(F\) 的不可约因子只能是各个 \(p_i\)。因此
\[
R^\times=\left\{a\prod_{i=1}^s p_i^{n_i}:a\in k^\times,\ n_i\in\mathbb Z\right\}.
\]
同样，\(1/h=F/D^n\) 给出 \(hF=D^n\)，故 \(h\mid D\)；反向已由 \(1/h=(D/h)/D\) 得到。

环 \(R\) 由 \(t,1/D\) 代数生成。又 \(k[t]\subseteq R\subseteq k(t)\)，取分式域得到 \(\operatorname{Frac}(R)=k(t)\)。然而 \(1+tD\) 是非常数多项式，且与 \(D\) 互素；取它的一个不可约因子 \(h\)，便有 \(1/h\notin R\)。所以 \(R\) 不是域。代数生成只允许对所选生成元作多项式运算，取分式域则允许对每个非零元素取逆，二者因此给出不同的对象。
\end{solution}

\begin{exercise}\label{b3:ex:06-finite-field-closed-points}
列出 \(\mathbb F_2[x]\) 中剩余域次数为一、二、三的全部极大理想。
\end{exercise}
\begin{solution}
一次不可约式为 \(x,x+1\)；二次为 \(x^2+x+1\)；三次为
\(x^3+x+1,x^3+x^2+1\)。二次或三次多项式可约时必有一次因子，所以只需排除在零或一取零的首一多项式，便得到完整清单。对应剩余域分别有 \(2,4,8\) 个元素。后两类不是 \(\mathbb F_2\) 有理点，但都是素谱中的闭点。
\end{solution}

\begin{exercise}\label{b3:ex:06-real-residue-field}
证明 \(\mathfrak m=(x^2+1,y-x)\) 是 \(\mathbb R[x,y]\) 的极大理想，并求其在 \(\mathbb C^2\) 中的零点。
\end{exercise}
\begin{solution}
先用 \(y=x\) 消元，商环为 \(\mathbb R[x]/(x^2+1)\cong\mathbb C\)，所以极大。它没有实零点；复数零点恰为 \((i,i)\) 与 \((-i,-i)\)。这说明一个基域上的闭点在扩域后可以分裂为多个坐标点。
\end{solution}

\begin{exercise}\label{b3:ex:06-unit-ideal-certificate}
证明 \(xy-1\) 与 \(y^2\) 在任意域上都没有公共零点，并给出单位理想的显式证据。
\end{exercise}
\begin{solution}
等式
\[
x^2y^2-(xy+1)(xy-1)=1
\]
成立，因此两者生成单位理想。若有公共零点，在该点取值将给出 \(0=1\)。此处已经给出证据，所以不需代数闭域假设。
\end{solution}

\begin{exercise}\label{b3:ex:06-empty-real-zero-set}
证明 \((x^2+y^2+1)\subseteq\mathbb R[x,y]\) 是真理想，却没有实零点，并给一个复零点。
\end{exercise}
\begin{solution}
非零的非常数多项式不是 \(\mathbb R[x,y]\) 的单位，因为非零乘积的全次数相加，所以所列主理想是真理想。实数上 \(x^2+y^2+1>0\)，无零点；复点 \((i,0)\) 满足方程。不能由实零点集为空推出单位理想。
\end{solution}

\begin{exercise}\label{b3:ex:06-fat-point-radical}
设 \(k\) 为域，令 \(I=(x^2,xy,y^3)\subseteq k[x,y]\)。计算 \(\sqrt I\) 和 \(k[x,y]/I\) 的一组 \(k\) 基；若 \(k\) 代数闭，说明点集遗漏了什么信息。
\end{exercise}
\begin{solution}
因为 \(x^2,y^3\in I\)，有 \((x,y)\subseteq\sqrt I\)。又 \(I\subseteq(x,y)\)，后者极大且为根理想，所以 \(\sqrt I=(x,y)\)。模掉 \(I\) 后只剩单项式 \(1,x,y,y^2\)；理想由单项式生成，这四项中没有任何非零线性组合属于 \(I\)，故为基。若 \(k\) 代数闭，普通零点集是 \(\{(0,0)\}\)，其坐标环为 \(k\)；原商环 \(k[x,y]/I\) 则有四维，且含有非零幂零元，这些信息不能由普通点集恢复。
\end{solution}

\begin{exercise}\label{b3:ex:06-hyperbola-vanishing}
设 \(k\) 代数闭，\(X=Z_k(xy-1)\)。证明 \(I(X)=(xy-1)\)，并计算 \(k[X]\)。
\end{exercise}
\begin{solution}
映射 \(k[x,y]\to k[t,t^{-1}]\)、\(x\mapsto t,y\mapsto t^{-1}\) 诱导商环同构：反向映射把 \(t\) 送到 \([x]\)，其逆元为 \([y]\)。商环为整环，所以 \((xy-1)\) 是素理想，因而为根理想。\cref{b3:thm:strong-nullstellensatz}于是给出 \(I(X)=(xy-1)\)，坐标环即 \(k[t,t^{-1}]\)。
\end{solution}

\begin{exercise}\label{b3:ex:06-unions-intersections}
对代数闭域上的代数集 \(X,Y\subseteq k^n\)，证明
\[
I(X\cup Y)=I(X)\cap I(Y),\qquad
I(X\cap Y)=\sqrt{I(X)+I(Y)}.
\]
\end{exercise}
\begin{solution}
多项式在并集上恒为零当且仅当在两集合上分别为零，给出第一式。第二式中，和理想的公共零点恰为 \(X\cap Y\)，因为 \(Z_k\bigl(I(X)\bigr)=X\)、\(Z_k\bigl(I(Y)\bigr)=Y\)；再应用\cref{b3:thm:strong-nullstellensatz}得到其零点理想为该和的根。
\end{solution}

\begin{exercise}\label{b3:ex:06-squaring-pullback}
设 \(k\) 为代数闭域。考察点集映射 \(k\to k\)、\(a\mapsto a^2\)，以及它诱导的坐标环同态 \(k[x]\to k[t]\)、\(x\mapsto t^2\)。证明点集映射满射，坐标环同态单射但不是满射；在特征二中，点集映射甚至双射。
\end{exercise}
\begin{solution}
代数闭性保证每个数有平方根，故点集映射满射。环同态为 \(P(x)\mapsto P(t^2)\)，不同次数产生不同的偶次幂，故单射；像只含偶次幂，所以不含 \(t\)，不是满射。特征二时 \(a^2=b^2\) 推出 \((a-b)^2=0\)，域中于是 \(a=b\)。因此点集双射也不能替代坐标环同构。
\end{solution}

\begin{exercise}\label{b3:ex:06-rabinowitsch-equivalence}
设 \(k\) 为任意域，\(R=k[x,y]\)、\(I=(x^2,xy,y^3)\)。在 \(R[T]\) 中令
\[
J_x=IR[T]+(1-Tx),\qquad J_{1+x}=IR[T]+(1-T(1+x)).
\]
给出 \(1\in J_x\) 的显式多项式组合；再证明 \(J_{1+x}\) 是真理想，并说明在 \(R/I\) 中倒置 \(x\) 与倒置 \(1+x\) 分别发生了什么。
\end{exercise}
\begin{solution}
因为 \(x^2\in I\)，有
\[
1=(1-Tx)(1+Tx)+T^2x^2,
\]
这就是第一项所需的单位理想证据。另一方面，取值 \((x,y,T)=(0,0,1)\) 使 \(J_{1+x}\) 的所有生成元为零，因此它是真理想，这一取值在任意域中都可作出。

在 \(C=R/I\) 中，\(\bar x^2=0\)，故将 \(\bar x\) 变成单位会迫使 \(C_{\bar x}=0\)。而 \(1+\bar x\) 已是单位，其逆为 \(1-\bar x\)，所以 \(C_{1+\bar x}=C\ne0\)。两个增广理想分别记录了这两种倒置的结果；\cref{b3:lem:rabinowitsch-equivalence}把它们与 \(x\in\sqrt I\)、\(1+x\notin\sqrt I\) 对应起来。
\end{solution}

\begin{exercise}\label{b3:ex:06-finite-field-vanishing}
设 \(q\) 为素数幂，\(n\ge1\) 为整数。证明在 \(\mathbb F_q[x_1,\ldots,x_n]\) 中，
\[
I(\mathbb F_q^n)=(x_1^q-x_1,\ldots,x_n^q-x_n).
\]
\end{exercise}
\begin{solution}
右侧的生成元在每个点上为零，故包含于左侧。反向，依次用这些关于不同变量的首一多项式除法，得到余式 \(r\)，其中每个变量的次数都小于 \(q\)，且 \(r\) 仍在全部点上为零。对变量数归纳：\(n=1\) 时非零次数小于 \(q\) 的多项式不可能有 \(q\) 个根；一般情形固定前 \(n-1\) 个变量，将 \(r\) 视为最后一个变量的多项式，所有系数在这些固定值上为零，再由归纳假设得系数多项式为零。因此 \(r=0\)，原多项式属于右侧理想。
\end{solution}

\begin{exercise}\label{b3:ex:06-nonreduced-normalization}
设 \(k\) 为域。按本书约定求 \(A=k[t,\varepsilon]/(\varepsilon^2)\) 的正规化，并说明为何 \(A\) 到正规化的自然映射不是单射。
\end{exercise}
\begin{solution}
写元素为 \(P(t)+\varepsilon\cdot Q(t)\)。幂零元恰好为 \((\varepsilon)\)：若一个幂为零，模掉 \(\varepsilon\) 后 \(P(t)\) 在整环 \(k[t]\) 中幂零，只能为零；反向该理想平方为零。因此 \(A_{\mathrm{red}}=k[t]\)，后者已整闭，正规化还是 \(k[t]\)。自然映射的核为非零理想 \((\varepsilon)\)。
\end{solution}



\begin{exercise}\label{b3:ex:06-three-axes}
设 \(k\) 为域，\(A=k[x,y,z]/(xy,xz,yz)\)。求全分式环、正规化、正规化商模的 \(k\) 维数及导子理想。
\end{exercise}
\begin{solution}
把一个元素送到三条坐标轴上的限制，得到单射
\[
A\longrightarrow B=k[x]\times k[y]\times k[z].
\]
确切地说，\(A\) 的元素唯一写成
\(c+x\cdot f(x)+y\cdot g(y)+z\cdot h(z)\)，其像恰为三个常数项相等的多项式组。因此环约化，极小素理想是三条轴对应的 \((y,z),(x,z),(x,y)\)。\cref{b3:prop:reduced-total-quotient-product}给出全分式环
\(k(x)\times k(y)\times k(z)\)。

在 \(Q(A)\) 中添入三个坐标幂等元 \(e_1,e_2,e_3\)，它们满足 \(e_i^2-e_i=0\)，且生成 \(B\)，所以 \(B/A\) 整。\(B\) 的三个分量各自整闭，任何在 \(A\) 上整的分数组逐分量属于这些多项式环。因此正规化为 \(B\)。

商模由三常数项模去对角常数组成，维数为二。一个 \(a\in A\) 若要满足 \(aB\subseteq A\)，其公共常数必须为零，因为分别乘 \(e_i\) 后三个常数仍须相等；反过来常数全零时乘任意 \(B\) 元素仍满足要求。因此导子为 \((x,y,z)_A\)，在 \(B\) 中即 \(xk[x]\times yk[y]\times zk[z]\)。
\end{solution}

\begin{exercise}\label{b3:ex:06-normal-base-birational}
设 \(A\subseteq B\) 为整环的整双有理扩张，且 \(A\) 整闭。证明 \(A=B\)。说明此处不需有限性。
\end{exercise}
\begin{solution}
双有理使 \(B\) 嵌在 \(\operatorname{Frac}(A)\) 中，每个 \(b\in B\) 又在 \(A\) 上整。按整闭的定义，\(b\in A\)。因此原包含为等号，仅用了逐元素整性。
\end{solution}

\begin{exercise}\label{b3:ex:06-infinite-conductor}
设 \(K/k\) 为无限次代数扩域，\(A=k+xK[x]\subseteq B=K[x]\)，如\cref{b3:exm:nonfinite-normalization}所示。计算导子与 \(B/A\)，并说明导子非零不保证扩张有限。
\end{exercise}
\begin{solution}
正次数部分都在 \(A\)，所以 \(B/A\cong K/k\) 作为 \(k\) 向量空间，且 \(xK[x]\) 零化此商。若导子中的元素常数项 \(c\in k\) 非零，取 \(\lambda\in K\setminus k\)，则与 \(\lambda\) 相乘的常数项 \(c\lambda\notin k\)，不再属于 \(A\)，矛盾。因此导子恰为 \(xK[x]\)，非零；但正文已由无限剩余域维数证明扩张不有限。
\end{solution}

\begin{exercise}\label{b3:ex:06-three-four-cusp}
设 \(k\) 为域，令 \(A=k[t^3,t^4]\)。证明其正规化为 \(B=k[t]\)，计算 \(B/A\) 的维数与导子。
\end{exercise}
\begin{solution}
\(t=t^4/t^3\) 在分式域中，且 \(t^3\in A\)，所以 \(B=A+At+At^2\) 有限；\(B\) 整闭使其成为正规化。由三、四生成的非负整数包括 \(0,3,4\) 及全部不小于六的整数，遗漏 \(1,2,5\)，因此 \(B/A\) 的基为 \([t],[t^2],[t^5]\)，维数三。

显然 \(t^6B\subseteq A\)，所以它包含于导子。若 \(a\in A\) 有非零常数项，则 \(at\) 的一次项非零；若常数项为零而三次项非零，则 \(at^2\) 的五次项非零；若前两项为零而四次项非零，则 \(at\) 的五次项非零。这些情况都不在导子中。因此导子中元素的这三个系数都为零，而 \(A\) 中其余可能出现的幂次都至少为六，故这些元素都属于 \(t^6B\)。于是导子恰为 \(t^6B=(t^6,t^7,t^8)_A\)。
\end{solution}

\begin{exercise}\label{b3:ex:06-two-five-cusp}
设 \(k\) 为域。对 \(A=k[t^2,t^5]\subseteq B=k[t]\)，求正规化商模与导子，并比较正文尖点例子。
\end{exercise}
\begin{solution}
\(t=t^5/(t^2)^2\)，又满足 \(T^2-t^2=0\)，故正规化为 \(k[t]\)。可出现的幂次是零、全部正偶数和不小于五的奇数，遗漏一、三，因此 \(B/A\) 有基 \([t],[t^3]\)，维数二。全部不小于四的单项式都在 \(A\)，所以 \(t^4B\) 在导子中；若 \(a\) 的常数项或二次项非零，乘 \(t\) 后会出现一次项或三次项，故不在导子中。因此导子为 \(t^4B=(t^4,t^5)_A\)。与 \(k[t^2,t^3]\) 相比，缺失的函数和导子的消失阶数都增加了。
\end{solution}

\begin{exercise}\label{b3:ex:06-trace-bound-quadratic}
在 \(\mathbb Q(\sqrt5)/\mathbb Q\) 中取由整元素组成的域基 \(u_1=1,u_2=\sqrt5\)。计算迹矩阵 \(T=\bigl(\operatorname{Tr}(u_iu_j)\bigr)\)，并用\cref{b3:lem:normalization-separable-finite}说明所有代数整数的公共分母受到怎样的控制。
\end{exercise}
\begin{solution}
迹把 \(a+b\sqrt5\) 送到 \(2a\)，故
\(T=\operatorname{diag}(2,10)\)、\(\det T=20\)。
若 \(z=a+b\sqrt5\) 在 \(\mathbb Z\) 上整，则 \(\operatorname{Tr}(z)=2a\) 与 \(\operatorname{Tr}(z\sqrt5)=10b\) 都是整数，故
\[
z\in \tfrac12\mathbb Z+\tfrac1{10}\mathbb Z\sqrt5,
\qquad
10z\in\mathbb Z+\mathbb Z\sqrt5.
\]
因此 \(10\) 已是统一公共分母，当然 \(20=\det T\) 也可以。伴随矩阵方法只需给出一个统一非零分母，并不保证所得分母最优。这给出有限格内的控制，不声称所有该格中的元素都整。例如 \((1+\sqrt5)/2\) 整，而 \(1/2\) 不整。
\end{solution}

\begin{exercise}\label{b3:ex:06-inseparable-finite-coefficients}
设 \(k\) 为特征 \(p>0\) 的域，\(a\in k\setminus k^p\)，\(x\) 为变量。令 \(K=k(x)\)，在 \(K\) 的一个代数闭包中取 \(z\) 满足 \(z^p=ax^p+x\)。在添入 \(\alpha=a^{1/p}\)、\(y=x^{1/p}\) 后，说明 \(z\) 如何消除，并给出 \(k[\alpha,y]\) 作为 \(k[x]\) 模的一组有限生成元。
\end{exercise}
\begin{solution}
在共同代数闭包中，
\[
(\alpha y^p+y)^p=ay^{p^2}+y^p=ax^p+x.
\]
特征 \(p\) 的域中取 \(p\) 次幂单射，所以 \(z=\alpha y^p+y\)。因而加入有限个系数根与参数根后，这个纯不可分生成元已属于新的有理函数域。环 \(k[\alpha,y]\) 在 \(k[x]\) 上由
\(\alpha^i y^j\)、\(0\le i,j<p\) 生成，因为 \(\alpha^p=a\)、\(y^p=x\)。本题不要求确定原环在 \(K(z)\) 中的整闭包；它展示的是有限基变换如何处理不可分部分。
\end{solution}

\begin{exercise}\label{b3:ex:06-inseparable-trace-zero}
设 \(k\) 为特征 \(p>0\) 的域，\(t\) 为变量，\(L=k(t)\)、\(K=k(t^p)\)。证明 \(L/K\) 次数为 \(p\)，且迹映射恒为零。为什么不能把可分情形的迹分母论证原样用于此例？
\end{exercise}
\begin{solution}
\(1,t,\ldots,t^{p-1}\) 在 \(K\) 上线性无关：清分母后，不同幂次模 \(p\) 的余数不同，非零项不能抵消；又 \(t^p\in K\)，故它们是基。对 \(1\le i<p\)，乘 \(t^i\) 在这组基上的矩阵没有对角项，因为指数的余数都被非零的 \(i\) 改变；乘常数 \(c\in K\) 的迹为 \(pc=0\)。线性性给出全部元素的迹为零。因此对任一 \(K\) 基 \(u_1,\ldots,u_p\)，迹矩阵 \(T=(\operatorname{Tr}_{L/K}(u_i u_j))\) 为零矩阵，\(\Delta=\det T=0\)。可分证明中用来限制分母的 \(\Delta^{-1}\) 在这里不存在。
\end{solution}

\begin{exercise}\label{b3:ex:06-conductor-local-support}
设 \(k\) 为域，\(A=k[t]\subseteq B=k[t,t^{-1}]\)，令 \(M=B/A\)。求导子 \(\mathfrak c=\operatorname{Ann}_A(M)\) 与 \(\operatorname{Supp}_A(M)\)，并证明 \(B\) 不是有限生成的 \(A\) 模。比较 \(\operatorname{Supp}_A(M)\) 与 \(V_A(\mathfrak c)\)，说明有限扩张下的导子局部同构判据为什么不能无条件沿用。
\end{exercise}
\begin{solution}
若 \(a\in\mathfrak c\)，则对每个 \(n\ge1\) 都有 \(at^{-n}\in k[t]\)，所以 \(a\) 被任意高次的 \(t\) 整除，只能为零。因此 \(\mathfrak c=0\)。有限多个 Laurent 多项式有共同的最低幂次下界，它们的 \(A\) 线性组合也有这个下界，不可能生成全部 \(t^{-n}\)，故 \(B\) 不是有限生成的 \(A\) 模。

当 \(t\notin\mathfrak p\) 时，\(t\) 在 \(A_{\mathfrak p}\) 中已经可逆，故 \(B\otimes_AA_{\mathfrak p}=A_{\mathfrak p}\)，即 \(M_{\mathfrak p}=0\)。在 \(\mathfrak p=(t)\) 处，类 \(t^{-1}+A\) 的局部化却非零：若它被某个 \(s\in A\setminus(t)\) 零化，就有 \(st^{-1}\in A\)，迫使 \(t\mid s\)，矛盾。因此
\[
\operatorname{Supp}_A(M)=\{(t)\}=V_A(t)
\quad\subsetneq\quad V_A(\mathfrak c)=\operatorname{Spec}A.
\]
例如在 \(\mathfrak p=(0)\) 处两环局部化相同，但 \(\mathfrak c\setminus\mathfrak p\) 为空。每个 \(M\) 中的元素都能被某个 \(t\) 的幂零化，却没有一个共同的非零元素零化整个 \(M\)；有限生成假设正是将逐元素的零化条件化为有限个条件、再取共同乘积的依据。
\end{solution}

\begin{exercise}\label{b3:ex:06-node-local-idempotents}
设 \(k\) 为特征不等于二的域，\(A=k[t^2-1,t(t^2-1)]\subseteq B=k[t]\)，\(\mathfrak m=(t^2-1,t(t^2-1))_A\)，\(S=A\setminus\mathfrak m\)。解释为什么节点的局部正规化 \(S^{-1}B\) 没有非平凡幂等元，而它模掉 \(\mathfrak mA_{\mathfrak m}\) 后的纤维有非平凡幂等元，并写出后者。
\end{exercise}
\begin{solution}
局部正规化是整环，若 \(e^2=e\)，则 \(e(e-1)=0\)，所以 \(e=0\) 或一。纤维为 \(k[t]/(t^2-1)\)，特征不为二时，
\[
e_+=\frac{t+1}{2},\qquad e_-=\frac{1-t}{2}
\]
满足 \(e_\pm^2=e_\pm\)、\(e_++e_-=1\)、\(e_+e_-=0\)，分别在 \(t=1,-1\) 取值 \((1,0)\)、\((0,1)\)。这两个幂等元是商环中的类，不在整环中满足同一恒等式。
\end{solution}

\begin{exercise}\label{b3:ex:06-node-characteristic-two}
设 \(k\) 为特征二的域，\(A=k\bigl[t^2-1,t(t^2-1)\bigr]\)。用新变量 \(u=t-1\) 把 \(A\) 化成尖点形式，并计算其导子和原点纤维。
\end{exercise}
\begin{solution}
此时 \(t^2-1=u^2\)，且 \(t(t^2-1)=u^3+u^2\)，所以
\(A=k[u^2,u^3]\)。正规化是 \(k[u]\)，由\cref{b3:prop:cusp-conductor}导子为 \(u^2k[u]\)。原点理想扩张为 \((u^2)\)，纤维为 \(k[u]/(u^2)\)，只有一个点且非约化。于是正文节点的“两点”结论确实需要特征不为二。
\end{solution}

\begin{exercise}\label{b3:ex:06-finite-type-normalization-algebra}
设 \(k\) 为域，\(A\) 是约化有限型 \(k\) 代数。证明其正规化 \(\widetilde A\) 仍为有限型 \(k\) 代数。再说明若 \(A\) 为整环，正规化不改变分式域的超越次数。
\end{exercise}
\begin{solution}
由\cref{b3:thm:normalization-finite}，取有限个 \(A\) 模生成元 \(b_1,\ldots,b_s\) 生成 \(\widetilde A\)。若 \(a_1,\ldots,a_r\) 作为 \(k\) 代数生成 \(A\)，则
\(a_1,\ldots,a_r,b_1,\ldots,b_s\) 作为 \(k\) 代数生成 \(\widetilde A\)：每个元素都是 \(b_j\) 的 \(A\) 线性组合。整环情形
\(A\subseteq\widetilde A\subseteq\operatorname{Frac}(A)\)，且两侧取分式域得到同一域，因此超越次数不变。
\end{solution}
